\section{基于集成的分类}
\label{sec:classification-ensemble}

第\ref{sec:classic-overview-ensemble}节给出成员组合的基本动机，第\ref{chap:ensemble-learning}章系统讨论Bagging、Boosting、Stacking及成员简化。本节只比较组合后交付什么分类信息，以及分类损失、阈值和概率校准怎样影响使用。回归树也可以作为分类集成的成员：GBDT用树拟合类别得分的修正，整体模型再给出类别决策；成员本身的输出类型不决定整体任务。

\begin{table*}[htbp]
  \centering
  \small
  \begin{tabularx}{\linewidth}{@{}>{\RaggedRight\arraybackslash}p{22mm}>{\RaggedRight\arraybackslash}X>{\RaggedRight\arraybackslash}X@{}}
    \toprule
    方法 & 分类输出 & 任务中的条件与代价 \\
    \midrule
    Bagging & 硬标签投票，或平均成员得分／概率后决策 & 概率平均不自动校准；预测成本随所用成员增加 \\
    随机森林 & 树投票或叶类别比例平均 & 树深、特征抽样与树数需选择；袋外诊断不替代测试 \\
    AdaBoost & 对弱分类器加权投票或输出累计得分 & 加权分类问题上的成员表现与噪声影响轮次选择 \\
    GBDT & 累加类别得分，再按阈值或Softmax决策 & 分类损失决定修正；树深、学习率与早停共同影响结果 \\
    Stacking & 元学习器组合标签、得分或概率 & 组合层用留出或折外基预测训练；校准仍需独立检验 \\
    \bottomrule
  \end{tabularx}
  \caption{集成方法的分类输出与使用条件。共同训练机制已在简介中说明；此处按相同分类目标比较，不把成员数量视为性能保证。}
  \label{tab:classification-ensemble-map}
\end{table*}

第\ref{sec:boosting}节的投票、弱学习与提升风险解释组合可能带来什么收益。实际评价仍须按表\ref{tab:classification-evaluation-protocol}区分标签错误、排序、概率质量与误判代价，并计入训练和预测所需的全部成员；单个成员的概率或解释性质不能直接当作整体性质。
